\chapter{RESULTADOS E DISCUSSÕES}

\section{Modelo de Aprendizado de Máquina}

O melhor modelo de AM desenvolvido neste estudo apresentou acurácia máxima de aproximadamente $60\%$ na tarefa de prever a pontuação final de um time em sua liga nacional, utilizando
como base os dados estatísticos de desempenho individual de seus jogadores na temporada anterior.

Para avaliar o desempenho dos diferentes algoritmos aplicados à tarefa de previsão, foram calculadas métricas 
amplamente utilizadas em problemas de regressão, incluindo o coeficiente de determinação (R\textsuperscript{2}), 
o erro absoluto médio (MAE), o erro quadrático médio (MSE), a raiz do erro quadrático médio (RMSE), 
o erro percentual absoluto médio (MAPE) e o tempo de processamento de cada modelo. A Tabela~\ref{tab:comparacao_modelos} 
apresenta um resumo comparativo dos resultados obtidos, permitindo analisar tanto a qualidade preditiva 
quanto o custo computacional de cada abordagem.


\begin{table}[H]
\centering
\caption{Comparação de desempenho entre os modelos de Machine Learning avaliados}
\label{tab:comparacao_modelos}
\begin{tabular}{lcccccc}
\hline
\textbf{Modelo} & \textbf{R\textsuperscript{2}} & \textbf{MAE} & \textbf{MSE} & \textbf{RMSE} & \textbf{MAPE} & \textbf{Tempo (s)} \\
\hline
RL                & 0.41 & 11.48 & 191.78 & 13.85 & 0.2580 & 2.15 \\
RF                & 0.28 & 12.05 & 234.05 & 15.30 & 0.2690 & 423.13 \\
XGBoost           & 0.41 & 11.48 & 191.78 & 13.85 & 0.2580 & 2.86 \\
MLP               & 0.48 & 11.08 & 170.42 & 13.05 & 0.2424 & 161.68 \\
\hline
\end{tabular}
\begin{tablenotes}
\small
\item Fonte: Elaborado pelo autor.
\end{tablenotes}
\end{table}




\section{Otimização de Montagem de Elenco}